\scanpage{85}
\mspage{78}
% archive sheet number: 84
% Three long diagonal cancellation strokes cover the written draft.
\deleted{
J'ai \unsure{envie} ici de préciser la forme des systèmes inductifs intervenant dans c, quitte à faire donc une hyp. supplémentaire. Je vais considérer, \inserted{$\pi\geq\pi_0$ fixé, \unsure{et|tt}} pour \deleted{\unsure{$X$ fixé}} tt $X$, l'ens. ordonné $I_\pi(X)$ défini ainsi
\[
\begin{aligned}
I_\pi(X)={}&\text{ens. des sous-objets }\inserted{X_i}\text{ de }X\text{ tels}\\
&\text{que }X_i\in\mathrm{Filt}^{\pi}(E),\text{ ordonné par}\\
&\text{inclusion.}
\end{aligned}
\]
Il faut voir que cet ensemble \inserted{ordonné} est grand devant $\pi$, et que
\[
\varinjlim_{i\in I_\pi(X)}X_i=X.
\]
Si $X\in\mathrm{Ob}\,\mathrm{Filt}^{\pi}(E)$, alors $X$ est un plus grand \deleted{objet} élément de $I_\pi(X)$, qui est donc grand devant \unsure{tt|tout}, et $\varinjlim X_i=X_X=X$.

Considérons \deleted{\unsure{donc un|donc|le cas}} le cas général $X\in\mathrm{Ob}\,\mathrm{Filt}^{\pi'}(E)$, avec $\pi'>\pi$. Soit \deleted{\unsure{un|Soit}} $J\subset I_\pi(X)$ avec $J$ de cardinal $\leq\pi$. Il faut donc voir qu'il existe $Y\subset X$, \deleted{\unsure{avec}} avec $Y\in\mathrm{Filt}^{\pi}_{\deleted{\unsure{\pi|0}}}(E)$, et $Y\supset X_i$ pour tout $i\in J$.
}
